\clearpage\selectlanguage{english}\fancyhead[L]{\small DISQUISITIONES GENERALES}
\SourceLeaf{134}
\begin{center}
DISQUISITIONES GENERALES

\medskip
CIRCA SERIEM INFINITAM
\[
\begin{aligned}&1+\frac{\alpha\beta}{1\mathbin{.}\gamma}x\\
&{}+\frac{\alpha(\alpha+1)\beta(\beta+1)}{1\mathbin{.}2\mathbin{.}\gamma(\gamma+1)}xx\\
&{}+\frac{\alpha(\alpha+1)(\alpha+2)\beta(\beta+1)(\beta+2)}{1\mathbin{.}2\mathbin{.}3\mathbin{.}\gamma(\gamma+1)(\gamma+2)}x^3+\mathrm{etc.}
\end{aligned}
\]
PARS PRIOR.

\medskip
\rule{0.13\linewidth}{0.4pt}

\medskip
\textit{INTRODUCTIO.}

\end{center}
\SixthArticle{1}
Series, quam in hac commentatione perscrutari suscipimus, tamquam functio
quatuor quantitatum $\alpha,\beta,\gamma,x$ spectari potest, quas ipsius \textit{elementa} vocabimus,
ordine suo elementum primum $\alpha$, secundum $\beta$, tertium $\gamma$, quartum $x$ distinguentes. Manifesto elementum primum cum secundo permutare licet: quodsi itaque brevitatis caussa seriem nostram hoc signo $F(\alpha,\beta,\gamma,x)$ denotamus, habebimus $F(\beta,\alpha,\gamma,x)=F(\alpha,\beta,\gamma,x)$.

\SixthArticle{2}
Tribuendo elementis $\alpha,\beta,\gamma$ valores determinatos, series nostra in functionem unicae variabilis $x$ transit, quae manifesto post terminum $1-\alpha^{\mathrm{tum}}$ vel
$1-\beta^{\mathrm{tum}}$ abrumpitur, si $\alpha-1$ vel $\beta-1$ est numerus integer negativus, in casibus reliquis vero in infinitum excurrit. In casu priori series exhibet functionem
algebraicam rationalem, in posteriori autem plerumque functionem transcendentem. Elementum tertium $\gamma$ debet esse neque numerus negativus integer neque
$=0$, ne ad terminos infinite magnos delabamur.
\SourceLeaf{135}
\SixthArticle{3}
Co\"efficientes potestatum $x^m$, $x^{m+1}$ in serie nostra sunt ut
\[
1+\frac{\gamma+1}{m}+\frac{\gamma}{mm}:1+\frac{\alpha+\beta}{m}+\frac{\alpha\beta}{mm}
\]
adeoque ad rationem aequalitatis eo magis accedunt, quo maior assumitur $m$. Si
itaque etiam elemento quarto $x$ valor determinatus tribuitur, ab huius indole
convergentia seu divergentia pendebit. Quoties scilicet ipsi $x$ tribuitur valor realis, positivus seu negativus, unitate minor, series certo, si non statim ab initio,
tamen post certum intervallum, convergens erit, atque ad summam finitam ex
asse determinatam perducet. Idem eveniet per valorem imaginarium ipsius $x$
formae $a+b\BareRoot{-1}$, quoties $aa+bb<1$. Contra pro valore ipsius $x$ reali unitateque maiori, vel pro imaginario formae $a+b\BareRoot{-1}$, quoties $aa+bb>1$, series si non statim tamen post certum intervallum necessario divergens erit, ita ut
de ipsius \textit{summa} sermo esse nequeat. Denique pro valore $x=1$ (seu generalius
pro valore formae $a+b\BareRoot{-1}$, quoties $aa+bb=1$) seriei convergentia seu divergentia ab ipsarum $\alpha,\beta,\gamma$ indole pendebit, de qua, atque in specie de summa
seriei pro $x=1$, in Sect. tertia loquemur.

Patet itaque, quatenus functio nostra tamquam summa seriei definita sit,
disquisitionem natura sua restrictam esse ad casus eos, ubi series revera convergat, adeoque quaestionem ineptam esse, quinam sit valor seriei pro valore ipsius
$x$ unitate maiori. Infra autem, inde a Sectione quarta, functionem nostram altiori principio superstruemus, quod applicationem generalissimam patiatur.

\SixthArticle{4}
Differentiatio seriei nostrae, considerando solum elementum quartum $x$
tamquam variabile, ad functionem similem perducit, quum manifesto habeatur
\[
\frac{\mathrm{d}F(\alpha,\beta,\gamma,x)}{\mathrm{d}x}
=\frac{\alpha\beta}{\gamma}F(\alpha+1,\beta+1,\gamma+1,x)
\]
Idem valet de differentiationibus repetitis.

\SixthArticle{5}
Operae pretium erit, quasdam functiones, quas ad seriem nostram reducere
licet, quarumque usus in tota analysi est frequentissimus, hic apponere.
\SourceLeaf{136}
\[
\begin{array}{rl}
\mathrm{I.}&(t+u)^n=t^nF(-n,\beta,\beta,-\frac{u}{t})
\end{array}
\]
ubi elementum $\beta$ est arbitrarium.
\[
\begin{array}{rl}
\mathrm{II.}&(t+u)^n+(t-u)^n=2t^nF(-\frac{1}{2}n,-\frac{1}{2}n+\frac{1}{2},\frac{1}{2},\frac{uu}{tt})
\end{array}
\]
\[
\begin{array}{rl}
\mathrm{III.}&(t+u)^n+t^n=2t^nF(-n,\omega,2\omega,-\frac{u}{t})
\end{array}
\]
denotante $\omega$ quantitatem infinite parvam.
\[
\begin{array}{rl}
\mathrm{IV.}&(t+u)^n-(t-u)^n=2nt^{n-1}uF(-\frac{1}{2}n+\frac{1}{2},-\frac{1}{2}n+1,\frac{3}{2},\frac{uu}{tt})
\end{array}
\]
\[
\begin{array}{rl}
\mathrm{V.}&(t+u)^n-t^n=nt^{n-1}uF(1-n,1.\ 2,-\frac{u}{t})
\end{array}
\]
\[
\begin{array}{rl}
\mathrm{VI.}&\log(1+t)=tF(1,1,2,-t)
\end{array}
\]
\[
\begin{array}{rl}
\mathrm{VII.}&\log\frac{1+t}{1-t}=2tF(\frac{1}{2},1,\frac{3}{2},tt)
\end{array}
\]
\[
\begin{array}{rl}
\mathrm{VIII.}&e^t=F(1,k,1,\frac{t}{k})=1+tF(1,k,2,\frac{t}{k})=1+t+\frac{1}{2}ttF(1,k,3,\frac{t}{k})\ \mbox{etc.}
\end{array}
\]
denotante $e$ basin logarithmorum hyperbolicorum, $k$ numerum infinite magnum.
\[
\begin{array}{rl}
\mathrm{IX.}&e^t+e^{-t}=2F(k,k',\frac{1}{2},\frac{tt}{4kk'})
\end{array}
\]
denotantibus $k$, $k'$ numeros infinite magnos.
\[
\begin{array}{rl}
\mathrm{X.}&e^t-e^{-t}=2tF(k,k',\frac{3}{2},\frac{tt}{4kk'})
\end{array}
\]
\[
\begin{array}{rl}
\mathrm{XI.}&\sin t=tF(k,k',\frac{3}{2},-\frac{tt}{4kk'})
\end{array}
\]
\[
\begin{array}{rl}
\mathrm{XII.}&\cos t=F(k,k',\frac{1}{2},-\frac{tt}{4kk'})
\end{array}
\]
\[
\begin{array}{rl}
\mathrm{XIII.}&t=\sin t.\ F(\frac{1}{2},\frac{1}{2},\frac{3}{2},\sin t^2)
\end{array}
\]
\[
\begin{array}{rl}
\mathrm{XIV.}&t=\sin t.\ \cos t.\ F(1,1,\frac{3}{2},\sin t^2)
\end{array}
\]
\[
\begin{array}{rl}
\mathrm{XV.}&t=\mathrm{tang}\,t.\ F(\frac{1}{2},1,\frac{3}{2},-\mathrm{tang}\,t^2)
\end{array}
\]
\[
\begin{array}{rl}
\mathrm{XVI.}&\sin nt=n\sin t.\ F(\frac{1}{2}n+\frac{1}{2},-\frac{1}{2}n+\frac{1}{2},\frac{3}{2},\sin t^2)
\end{array}
\]
\[
\begin{array}{rl}
\mathrm{XVII.}&\sin nt=n\sin t.\ \cos t.\ F(\frac{1}{2}n+1,-\frac{1}{2}n+1,\frac{3}{2},\sin t^2)
\end{array}
\]
\[
\begin{array}{rl}
\mathrm{XVIII.}&\sin nt=n\sin t.\ \cos t^{n-1}F(-\frac{1}{2}n+1,-\frac{1}{2}n+\frac{1}{2},\frac{3}{2},-\mathrm{tang}\,t^2)
\end{array}
\]
\[
\begin{array}{rl}
\mathrm{XIX.}&\sin nt=n\sin t.\ \cos t^{-n-1}F(\frac{1}{2}n+1,\frac{1}{2}n+\frac{1}{2},\frac{3}{2},-\mathrm{tang}\,t^2)
\end{array}
\]
\[
\begin{array}{rl}
\mathrm{XX.}&\cos nt=F(\frac{1}{2}n,-\frac{1}{2}n,\frac{1}{2},\sin t^2)
\end{array}
\]
\[
\begin{array}{rl}
\mathrm{XXI.}&\cos nt=\cos t.\ F(\frac{1}{2}n+\frac{1}{2},-\frac{1}{2}n+\frac{1}{2},\frac{1}{2},\sin t^2)
\end{array}
\]
\[
\begin{array}{rl}
\mathrm{XXII.}&\cos nt=\cos t^nF(-\frac{1}{2}n,-\frac{1}{2}n+\frac{1}{2},\frac{1}{2},-\mathrm{tang}\,t^2)
\end{array}
\]
\[
\begin{array}{rl}
\mathrm{XXIII.}&\cos nt=\cos t^{-n}F(\frac{1}{2}n+\frac{1}{2},\frac{1}{2}n,\frac{1}{2},-\mathrm{tang}\,t^2)
\end{array}
\]
\SourceLeaf{137}
\SixthArticle{6}
Functiones praecedentes sunt algebraicae atque transcendentes a logarithmis circuloque pendentes. Neutiquam vero \textit{harum} caussa disquisitionem nostram
\textit{generalem} suscipimus, sed potius in gratiam theoriae functionum transcendentium altiorum promovendae, quarum genus amplissimum series nostra complectitur. Huc, inter infinita alia, pertinent co\"efficientes ex evolutione functionis
$(aa+bb-2ab\cos\varphi)^{-n}$ in seriem secundum cosinus angulorum $\varphi$, $2\varphi$, $3\varphi$ etc.
progredientem orti, de quibus \textit{in specie} alia occasione fusius agemus. Ad formam
seriei nostrae autem illi co\"efficientes pluribus modis reduci possunt. Scilicet statuendo
\[
(aa+bb-2ab\cos\varphi)^{-n}=\Omega=A+2A'\cos\varphi+2A''\cos2\varphi+2A'''\cos3\varphi+\mbox{etc.}
\]
habemus \textit{primo}
\[
\begin{array}{rl}
A&=a^{-2n}F(n,n,1,\frac{bb}{aa})
\end{array}
\]
\[
\begin{array}{rl}
A'&=na^{-2n-1}bF(n,n+1,2,\frac{bb}{aa})
\end{array}
\]
\[
\begin{array}{rl}
A''&=\frac{n(n+1)}{1\mathbin{.}2}a^{-2n-2}bbF(n,n+2,3,\frac{bb}{aa})
\end{array}
\]
\[
\begin{array}{rl}
A'''&=\frac{n(n+1)(n+2)}{1\mathbin{.}2\mathbin{.}3}a^{-2n-3}b^3F(n,n+3,4.\ \frac{bb}{aa})
\end{array}
\]
\begin{center}etc.\end{center}
Si enim $aa+bb-2ab\cos\varphi$ consideratur tamquam productum ex $a-br$ in
$a-br^{-1}$ (designante $r$ quantitatem $\cos\varphi+\sin\varphi.\BareRoot{-1}$), fit $\Omega$ aequalis producto
\[
\begin{array}{rl}
\mbox{ex}&a^{-2n}
\end{array}
\]
\[
\begin{array}{rl}
\mbox{in}&1+n\frac{br}{a}+\frac{n(n+1)}{1\mathbin{.}2}.\ \frac{bbrr}{aa}+\frac{n(n+1)(n+2)}{1\mathbin{.}2\mathbin{.}3}.\ \frac{b^3r^3}{a^3}+\mbox{etc.}
\end{array}
\]
\[
\begin{array}{rl}
\mbox{in}&1+n\frac{br^{-1}}{a}+\frac{n(n+1)}{1\mathbin{.}2}.\ \frac{bbr^{-2}}{aa}+\frac{n(n+1)(n+2)}{1\mathbin{.}2\mathbin{.}3}.\ \frac{b^3r^{-3}}{a^3}+\mbox{etc.}
\end{array}
\]
Quod productum quum identicum esse debeat cum
\[
A+A'(r+r^{-1})+A''(rr+r^{-2})+A'''(r^3+r^{-3})
\]
valores supra dati sponte prodeunt.

Porro habemus \textit{secundo}
\SourceLeaf{138}
\[
\begin{array}{rl}
A&=(aa+bb)^{-n}F(\frac12 n,\frac12 n+\frac12 .\,1,\frac{4aabb}{(aa+bb)^2})
\end{array}
\]
\[
\begin{array}{rl}
A'&=n(aa+bb)^{-n-1}abF(\frac12 n+\frac12,\frac12 n+1,2,\frac{4aabb}{(aa+bb)^2})
\end{array}
\]
\[
\begin{array}{rl}
A''&=\frac{n(n+1)}{1\mathbin{.}2}(aa+bb)^{-n-2}aabbF(\frac12 n+1,\frac12 n+\frac32,3,\frac{4aabb}{(aa+bb)^2})
\end{array}
\]
\[
\begin{array}{rl}
A'''&=\frac{n(n+1)(n+2)}{1\mathbin{.}2\mathbin{.}3}(aa+bb)^{-n-3}a^3b^3F(\frac12 n+\frac32,\frac12 n+2,4,\frac{4aabb}{(aa+bb)^2})
\end{array}
\]
\begin{center}etc.\end{center}
qui valores facile deducuntur ex
\[
\Omega(aa+bb)^n=1+n(r+r^{-1})\frac{ab}{aa+bb}+\frac{n(n+1)}{1\mathbin{.}2}(r+r^{-1})^2\frac{aabb}{(aa+bb)^2}+\mathrm{etc.}
\]
\textit{Tertio} fit
\[
\begin{array}{rl}
A&=(a+b)^{-2n}F(n,\frac12,1,\frac{4ab}{(a+b)^2})
\end{array}
\]
\[
\begin{array}{rl}
A'&=n(a+b)^{-2n-2}abF(n+1,\frac32,3,\frac{4ab}{(a+b)^2})
\end{array}
\]
\[
\begin{array}{rl}
A''&=\frac{n(n+1)}{1\mathbin{.}2}(a+b)^{-2n-4}aabbF(n+2,\frac52,5,\frac{4ab}{(a+b)^2})
\end{array}
\]
\[
\begin{array}{rl}
A'''&=\frac{n(n+1)(n+2)}{1\mathbin{.}2\mathbin{.}3}(a+b)^{-2n-6}a^3b^3F(n+3,\frac72,7,\frac{4ab}{(a+b)^2})
\end{array}
\]
\begin{center}etc.\end{center}
Denique fit \textit{quarto}
\[
\begin{array}{rl}
A&=(a-b)^{-2n}F(n,\frac12,1,-\frac{4ab}{(a-b)^2})
\end{array}
\]
\[
\begin{array}{rl}
A'&=n(a-b)^{-2n-2}abF(n+1,\frac32,3,-\frac{4ab}{(a-b)^2})
\end{array}
\]
\[
\begin{array}{rl}
A''&=\frac{n(n+1)}{1\mathbin{.}2}(a-b)^{-2n-4}aabbF(n+2,\frac52,5,-\frac{4ab}{(a-b)^2})
\end{array}
\]
\[
\begin{array}{rl}
A'''&=\frac{n(n+1)(n+2)}{1\mathbin{.}2\mathbin{.}3}(a-b)^{2n-6}a^3b^3F(n+3,\frac72,7,-\frac{4ab}{(a-b)^2})
\end{array}
\]
\begin{center}etc.\end{center}
Valores illi atque hi facile eruuntur ex
\[
\begin{array}{rl}
\Omega(a+b)^{2n}&=(1-\frac{4ab\cos\frac12\varphi^2}{(a+b)^2})^{-n}
\end{array}
\]
\[
\begin{array}{rl}
&=1+n\frac{ab}{(a+b)^2}(r^{\frac12}+r^{-\frac12})^2+\frac{n(n+1)}{1\mathbin{.}2}\mathbin{.}\frac{aabb}{(a+b)^4}(r^{\frac12}+r^{-\frac12})^4+\mathrm{etc.}
\end{array}
\]
\[
\begin{array}{rl}
\Omega(a-b)^{2n}&=(1+\frac{4ab\sin\frac12\varphi^2}{(a-b)^2})^{-n}
\end{array}
\]
\[
\begin{array}{rl}
&=1+n\frac{ab}{(a-b)^2}(r^{\frac12}-r^{-\frac12})^2+\frac{n(n+1)}{1\mathbin{.}2}\mathbin{.}\frac{aabb}{(a-b)^4}(r^{\frac12}-r^{-\frac12})^4+\mathrm{etc.}
\end{array}
\]
\begin{center}\rule{.23\linewidth}{.4pt}\end{center}
\SourceLeaf{139}
\begin{center}\textit{SECTIO PRIMA.}\\[4pt]
\textit{Relationes inter functiones contiguas:}\end{center}
\SixthArticle{7}
\textit{Functionem} ipsi $F(\alpha,\beta,\gamma,x)$ \textit{contiguam} vocamus, quae ex illa oritur, dum
elementum primum, secundum, vel tertium unitate vel augetur vel diminuitur,
manentibus tribus reliquis elementis. Functio itaque primaria $F(\alpha,\beta,\gamma,x)$ sex
contiguas suppeditat, inter quarum binas ipsamque primariam aequatio persimplex linearis datur. Has aequationes, numero quindecim, hic in conspectum
producimus, brevitatis gratia elementum quartum quod semper subintelligitur $=x$
omittentes, functionemque primariam simpliciter per $F$ denotantes.
\[
\begin{array}{rcl}
{[1]}&0=&(\gamma-2\alpha-(\beta-\alpha)x)F+\alpha(1-x)F(\alpha+1,\beta,\gamma)\\
&&{}-(\gamma-\alpha)F(\alpha-1,\beta,\gamma)
\end{array}
\]
\[
\begin{array}{rcl}
{[2]}&0=&(\beta-\alpha)F+\alpha F(\alpha+1,\beta,\gamma)-\beta F(\alpha,\beta+1,\gamma)
\end{array}
\]
\[
\begin{array}{rcl}
{[3]}&0=&(\gamma-\alpha-\beta)F+\alpha(1-x)F(\alpha+1,\beta,\gamma)-(\gamma-\beta)F(\alpha,\beta-1,\gamma)
\end{array}
\]
\[
\begin{array}{rcl}
{[4]}&0=&\gamma(\alpha-(\gamma-\beta)x)F-\alpha\gamma(1-x)F(\alpha+1,\beta,\gamma)\\
&&{}+(\gamma-\alpha)(\gamma-\beta)xF(\alpha,\beta,\gamma+1)
\end{array}
\]
\[
\begin{array}{rcl}
{[5]}&0=&(\gamma-\alpha-1)F+\alpha F(\alpha+1,\beta,\gamma)-(\gamma-1)F(\alpha,\beta,\gamma-1)
\end{array}
\]
\[
\begin{array}{rcl}
{[6]}&0=&(\gamma-\alpha-\beta)F-(\gamma-\alpha)F(\alpha-1,\beta,\gamma)+\beta(1-x)F(\alpha,\beta+1,\gamma)
\end{array}
\]
\[
\begin{array}{rcl}
{[7]}&0=&(\beta-\alpha)(1-x)F-(\gamma-\alpha)F(\alpha-1,\beta,\gamma)+(\gamma-\beta)F(\alpha,\beta-1,\gamma)
\end{array}
\]
\[
\begin{array}{rcl}
{[8]}&0=&\gamma(1-x)F-\gamma F(\alpha-1,\beta,\gamma)+(\gamma-\beta)xF(\alpha,\beta,\gamma+1)
\end{array}
\]
\[
\begin{array}{rcl}
{[9]}&0=&(\alpha-1-(\gamma-\beta-1)x)F+(\gamma-\alpha)F(\alpha-1,\beta,\gamma)\\
&&{}-(\gamma-1)(1-x)F(\alpha,\beta,\gamma-1)
\end{array}
\]
\[
\begin{array}{rcl}
{[10]}&0=&(\gamma-2\beta+(\beta-\alpha)x)F+\beta(1-x)F(\alpha,\beta+1,\gamma)\\
&&{}-(\gamma-\beta)F(\alpha,\beta-1,\gamma)
\end{array}
\]
\[
\begin{array}{rcl}
{[11]}&0=&\gamma(\beta-(\gamma-\alpha)x)F-\beta\gamma(1-x)F(\alpha,\beta+1,\gamma)\\
&&{}-(\gamma-\alpha)(\gamma-\beta)F(\alpha,\beta,\gamma+1)
\end{array}
\]
\[
\begin{array}{rcl}
{[12]}&0=&(\gamma-\beta-1)F+\beta F(\alpha,\beta+1,\gamma)-(\gamma-1)F(\alpha,\beta,\gamma-1)
\end{array}
\]
\[
\begin{array}{rcl}
{[13]}&0=&\gamma(1-x)F-\gamma F(\alpha,\beta-1,\gamma)+(\gamma-\alpha)xF(\alpha,\beta,\gamma+1)
\end{array}
\]
\[
\begin{array}{rcl}
{[14]}&0=&(\beta-1-(\gamma-\alpha-1)x)F+(\gamma-\beta)F(\alpha,\beta-1,\gamma)\\
&&{}-(\gamma-1)(1-x)F(\alpha,\beta,\gamma-1)
\end{array}
\]
\[
\begin{array}{rcl}
{[15]}&0=&\gamma(\gamma-1-(2\gamma-\alpha-\beta-1)x)F\\
&&{}+(\gamma-\alpha)(\gamma-\beta)xF(\alpha,\beta,\gamma+1)\\
&&{}-\gamma(\gamma-1)(1-x)F(\alpha,\beta,\gamma-1)
\end{array}
\]
\SixthArticle{8}
Ecce iam demonstrationem harum formularum. Statuendo
\[
\frac{(\alpha+1)(\alpha+2)\mathinner{.\,.\,.}(\alpha+m-1)\beta(\beta+1)\mathinner{.\,.\,.\,.\,.}(\beta+m-2)}{1\mathbin{.}2\mathbin{.}3\mathinner{.\,.\,.\,.\,.}m\mathbin{.}\gamma(\gamma+1)\mathinner{.\,.\,.\,.\,.}(\gamma+m-1)}=M
\]
\SourceLeaf{140}
erit. co\"efficiens potestatis $x^m$
\[
\begin{array}{lll}
\mbox{in }F&.\quad.\quad.\quad.\quad.\quad.&\alpha(\beta+m-1)M\\
\mbox{in }F(\alpha,\beta-1,\gamma)&.\quad.&\alpha(\beta-1)M\\
\mbox{in }F(\alpha+1,\beta,\gamma)&.\quad.&(\alpha+m)(\beta+m-1)M\\
\mbox{in }F(\alpha,\beta,\gamma-1)&.\quad.&\frac{\alpha(\beta+m-1)(\gamma+m-1)M}{\gamma-1}
\end{array}
\]
co\"efficiens autem potestatis $x^{m-1}$ in $F(\alpha+1,\beta,\gamma)$, seu co\"efficiens potestatis $x^m$
in $xF(\alpha+1,\beta,\gamma)$
\[
=m(\gamma+m-1)M
\]
Hinc statim demanat veritas formularum 5 et 3; permutando $\alpha$ cum $\beta$, oritur
ex 5 formula 12, atque ex his duabus per eliminationem 2. Perinde per eandem permutationem ex 3 oritur 6; ex 6 et 12 combinatis oritur 9, hinc per
permutationem 14, quibus combinatis habetur 7; denique ex 2 et 6 eruitur 1,
atque hinc permutando 10. Formula 8 simili modo ut supra formulae 5 et 3,
e consideratione co\"efficientium derivari potest (eodemque modo, si placeret, \textit{omnes}
15 formulae erui possent), vel elegantius ex iam notis sequenti modo. Mutando
in formula 5 elementum $\alpha$ in $\alpha-1$, atque $\gamma$ in $\gamma+1$, prodit
\[
0=(\gamma-\alpha+1)F(\alpha-1,\beta,\gamma+1)+(\alpha-1)F(\alpha,\beta,\gamma+1)-\gamma F(\alpha-1,\beta,\gamma)
\]
Mutando vero in formula 9 tantummodo $\gamma$ in $\gamma+1$, fit
\[
\begin{aligned}0={}&(\alpha-1-(\gamma-\beta)x)F(\alpha,\beta,\gamma+1)\\
&{}+(\gamma-\alpha+1)F(\alpha-1,\beta,\gamma+1)-\gamma(1-x)F(\alpha,\beta,\gamma)\end{aligned}
\]
E subtractione harum formularum statim oritur 8, atque hinc per permutationem 13. Ex 1 et 8 prodit 4, hincque permutando 11. Denique ex 8 et 9
deducitur 15.
\SixthArticle{9}
Si $\alpha'-\alpha$, $\beta'-\beta$, $\gamma'-\gamma$, nec non $\alpha''-\alpha$, $\beta''-\beta$, $\gamma''-\gamma$ sunt numeri integri (positivi seu negativi), a functione $F(\alpha,\beta,\gamma)$ ad functionem $F(\alpha',\beta',\gamma')$,
et perinde ab hac usque ad functionem $F(\alpha'',\beta'',\gamma'')$ transire licet per seriem similium functionum, ita ut quaelibet contigua sit antecedenti et consequenti, mutando scilicet primo elementum unum e.g. $\alpha$ continuo unitate, donec a $F(\alpha,\beta,\gamma)$
perventum sit ad $F(\alpha',\beta,\gamma)$, dein mutando elementum secundum, donec perven\SourceLeaf{141}%
tum sit ad $F(\alpha',\beta',\gamma)$, denique mutando elementum tertium, donec perventum sit
ad $F(\alpha',\beta',\gamma')$, et perinde ab hac usque ad $F(\alpha'',\beta'',\gamma'')$. Quum itaque per art. 7
habeantur aequationes lineares inter functionem primam, secundam atque tertiam,
et generaliter inter ternas quascunque consequentes huius seriei, facile perspicitur, hinc per eliminationem deduci posse aequationem linearem inter functiones
$F(\alpha,\beta,\gamma)$, $F(\alpha',\beta',\gamma')$, $F(\alpha'',\beta'',\gamma'')$, ita ut generaliter loquendo e duabus functionibus, quarum tria elementa prima numeris integris differunt, quamlibet aliam
functionem eadem proprietate gaudentem derivare liceat, siquidem elementum
quartum idem maneat. Ceterum hic nobis sufficit, hanc veritatem insignem generaliter stabilivisse, neque hic compendiis immoramur, per quae operationes ad
hunc finem necessariae quam brevissimae reddantur.
\SixthArticle{10}
Propositae sint e. g. functiones
\[
F(\alpha,\beta,\gamma),\ F(\alpha+1,\beta+1,\gamma+1),\ F(\alpha+2,\beta+2,\gamma+2)
\]
inter quas aequationem linearem invenire oporteat. Iungamus ipsas per functiones contiguas sequenti modo:
\[
\begin{array}{rl}
F(\alpha,\beta,\gamma)&=F
\end{array}
\]
\[
\begin{array}{rl}
F(\alpha+1,\beta,\gamma)&=F'
\end{array}
\]
\[
\begin{array}{rl}
F(\alpha+1,\beta+1,\gamma)&=F''
\end{array}
\]
\[
\begin{array}{rl}
F(\alpha+1,\beta+1,\gamma+1)&=F'''
\end{array}
\]
\[
\begin{array}{rl}
F(\alpha+2,\beta+1,\gamma+1)&=F''''
\end{array}
\]
\[
\begin{array}{rl}
F(\alpha+2,\beta+2,\gamma+1)&=F'''''
\end{array}
\]
\[
\begin{array}{rl}
F(\alpha+2,\beta+2,\gamma+2)&=F''''''
\end{array}
\]
Habemus itaque quinque aequationes lineares (e formulis 6, 13, 5 art. 7):
\[
\begin{array}{rcl}
\mathrm{I.}&0=&(\gamma-\alpha-1)F-(\gamma-\alpha-1-\beta)F'-\beta(1-x)F''
\end{array}
\]
\[
\begin{array}{rcl}
\mathrm{II.}&0=&\gamma F'-\gamma(1-x)F''-(\gamma-\alpha-1)xF'''
\end{array}
\]
\[
\begin{array}{rcl}
\mathrm{III.}&0=&\gamma F''-(\gamma-\alpha-1)F'''-(\alpha+1)F''''
\end{array}
\]
\[
\begin{array}{rcl}
\mathrm{IV.}&0=&(\gamma-\alpha-1)F'''-(\gamma-\alpha-2-\beta)F''''-(\beta+1)(1-x)F'''''
\end{array}
\]
\[
\begin{array}{rcl}
\mathrm{V.}&0=&(\gamma+1)F''''-(\gamma+1)(1-x)F'''''-(\gamma-\alpha-1)xF''''''
\end{array}
\]
Ex I et II prodit, eliminando $F'$

